\newpage
\phantomsection
\label{prf:common-refinement-partitions}
\LRAProofFor{lem:common-refinement-partitions}

\begin{remark*}[Return]
\hyperref[lem:common-refinement-partitions]{Return to Lemma}
\end{remark*}

\begin{lemma*}[Common Refinement of Two Partitions]
If \(P\) and \(P'\) are partitions of \([a,b]\), then \(P\cup P'\), ordered
in increasing order, is a partition of \([a,b]\) that refines both.
\end{lemma*}

\begin{proof}[Professional Standard Proof]
\LRAProofBodyStart
Both partitions are finite subsets of \([a,b]\) containing \(a\) and \(b\).
Their union is finite, contains \(a\) and \(b\), and becomes a strictly
increasing list after duplicate points are removed. Since \(P\subseteq P\cup
P'\) and \(P'\subseteq P\cup P'\), the ordered union refines both partitions.
\end{proof}

\begin{proof}[Detailed Learning Proof]
\LRAProofBodyStart
The only possible obstruction is ordering. A finite set of real numbers can be
listed increasingly, and duplicate common nodes are written once. The first
node is still \(a\), the last node is still \(b\), and every old node from both
partitions remains present.
\end{proof}

\begin{remark*}[Proof structure]
Order the finite union and observe that both original partitions are subsets
of it.
\end{remark*}

\begin{dependencies}
\begin{itemize}
  \item \hyperref[def:integration-partition]{Partition of a Compact Interval}.
\end{itemize}
\end{dependencies}

\clearpage
